\documentclass[11pt,letterpaper]{article}
\usepackage{amssymb}

\setlength{\parindent}{0pt}

\begin{document}

\begin{center}
{\Large \bf Numerical Linear Algebra }\\
{\Large \bf Project }\\
\end{center}

\small

Choose one of the three problems below. Do not do all three. You will only be
graded on one problem. The one you want me to grade.

Please ask me by email if directions below are not clear.

Codes need to be in ``julia''.\\
 
{\Large \bf Logistic}\\

Three topics to choose from: 
\begin{enumerate}
\item NEP --
nonsymmetric eigenvalue problem: Given an $n$-by-$n$ real nonsymmetric matrix, $A$,
compute ``quasi Schur form'': $A = Q T Q^T$.
\item SEP -- symmetric eigenvalue problem: Given an $n$-by-$n$ real symmetric matrix, $A$,
compute $A=QDQ^T$.
\item SVD -- singular value decomposition: Given an $m$-by-$n$ real matrix, $A$,
compute $A=USV^T$.
\end{enumerate}

Choose whichever one you prefer; more information on the topics is below. For
SVD, the book by Darve and Wooters is very well done (Section 5.4). In an ideal world,
you follow the book and you're good to go. For NEP, there's already a lot
done in ./src/gees.jl in Eric Darve's GitHub; obviously, that being said it's not just a matter
of copying and pasting the code and submitting it to me. Yes, you can reuse
everything. And I encourage you to look at everything carefully, but there's
extra work to do. For SEP, you need to ``symmetrize NEP.'' See Darve and
Wooters' book too. For SEP and SVD, the problems are intrinsically simpler to
solve, but you're starting from scratch. For NEP, you have some building blocks
in Darve's GitHub.\\

Please deliver as well some tests and some examples of useage of your code.  It
will be important that you assert that your code is correct and usable.  That
way you will also have to demonstrate your understanding of the properties of
the input and output of your code.

The projects must be done individually. You can communicate with each other,
but no copying code. I didn't have any problems with cheating on assignments 1,
2, and 3, so I'm not worried about the end-of-semester projects. Keep up the
good work.\\
 
Deliverable: The written report is simply your code. Either a .jl file or a
.ipynb file. (Or a mix.) No need to add a PDF, etc. I need to be able to execute your code. \\

The grade will be a "holistic" grade, so it's not 10\% for the oral
presentation and 40\% for the written exam, but it's a grade for everything.
Ultimately, what matters most are the codes, you need to be able to explain the
codes clearly during the oral presentation.\\

{\Large \bf Project \#1: NEP, nonsymmetric eigenvalue problem}\\

In input a nonsymmetric $n$-by-$n$ matrix $A$. In output the Schur form of $A$.
( $A = Q T Q^T$. ) Only real arithmetic, no complex arithmetic. $T$ is quasi
triangular so 1x1 block for real eigenvalues and 2x2 block for complex
eigenvalues. Please have a look at gees.jl in ./src of GitHub of Eric Darve for
a template. Algorithms needs to be the implicitly shited QR iteration algorithm using
Francis shift. Please first reduce to Hessenberg form, then each iteration
needs to be $\mathcal{O}(n^2)$ and so needs to take into account the Hessenberg structure of $H$.
Three things not done in gees.jl that you need to do (1/3) Need flags for
either 'D' compute the diagonal only, 'T' compute the Schur form only, or 'V'
compute both $Q$ and $T$. (2/3) Need to fix the ``explicit shift strategy''.
(3/3) Return the 2x2 block (whenever they exist on the diagonal) as ``standard''
form. See DM3 question 3.\\

Goal is that your code is able to work on a random matrix. Then you can try on a
matrix of the type $A = VDV^{-1}$ when $V$ is ill conditioned. Also you can try with cluster of
eigenvalues. A few other types of matrices are interesting.\\

Analysis: (*) please provide an example where we observe quadratic convergence for bottom right side of matrices,
(*) return a metric such as ``number of QR iterations per eigenvalue'', (*) time
your code for various $n$ and try to scale up to larger $n$ (say up to size $300$), 
compare with native julia functionality.\\

Rubrique: your code works for random matrices, is clearly written
and the analysis (stability, convergence and timing) is done, the choices made
in the design of the code are clearly explained in the comments of the code, you are able to clearly 
explain the code and answer questions 
during the oral presentation.\\

% How to get more points:

% (*) From the Schur form, write the computation for eigenvectors. Check that you
% have $A = V D V^{-1}$ (if $A$ is diagonalizable ``enough'').

% (*) Reorder the eigenvalues in the Schur form (and that associated Schur
% vector) such that all eigenvalues with positive real parts are at the start of
% the Schur form, and all the eigenvalues with negative real parts are at the
% end. 

% (*) Explore EAD: Early Aggressive Deflation

{\Large \bf Project \#2: SEP, symmetric eigenvalue problem}\\

In input a symmetric $n$-by-$n$ matrix $A$. ($A=A^T$.) In output the
eigendecomposition of $A$. ( $A = Q D Q^T$. ) $D$ is diagonal real and $Q$ is
$n$-by-$n$ orthogonal: $Q^TQ = QQ^T = I$. Only real arithmetic, no complex
arithmetic. (So do not take $A$ as complex Hermitian ($A=A^H$).) Real symmetric
is enough. First step is to take $A$ and reduce to symmetric tridiagonal form.
You can only use the lower part of $A$. (Or the upper, as you wish. But only
one half.) I recommend starting with gehrd! and then adapting for the symmetric
case. But as you wish. The nonsymmetric QR algorithm that is explained in Darve and Wootters 
needs to be ``symmetrized''. This is not that easy. Your final algorithm should only work on 
one half of a matrix, use $n^2/2+\mathcal{O}(n)$ memory, and perform $4/3 n^3+\mathcal{O}(n^2)$ FLOPS.\\

Need flags for
either ('D') compute the diagonal only or 'V'
compute both $Q$ and $D$. 
The goal is to perform $\mathcal{O}(n)$ operations per QR step if eigenvalues only are requested,
and $\mathcal{O}(n^2)$ operations per QR step if eigenvalues and eigenvectors are requested.\\

Analysis: (*) please provide an example where we observe cubic convergence for bottom right side of matrices,
(*) return a metric such as ``number of QR iterations per eigenvalue'', (*) time
your code for various $n$ and try to scale up to larger $n$.
(say up to size $300$), 
compare with native julia functionality.\\

Rubrique: your code works for random matrices, is clearly written
and the analysis (stability, convergence and timing) is done, the choice made
in the design of the code are clearly explained in the comments of the code, you are able to clearly 
explain the code and answer questions 
during the oral presentation.\\

% How to get more points:

% (*) Write the code for complex Hermitian matrix $A$ in input. The code needs to reduce
% $A$ (which is complex Hermitian) to a real symmetric tridiagonal matrix.
% Then you can reuse your previous code for a real symmetric tridiagonal matrices.

% (*) Explore EAD: Early Aggressive Deflation

% (*) Explore Divide and Conquer method.

{\Large \bf Project \#3: SVD, singular value decomposition}\\

In input an $m$-by-$n$ matrix $A$. (Please take $m\geq n$.) In output the thin SVD
of $A$. ( $A = U S V^T$. ) Only real arithmetic, no complex arithmetic. Please
read section 5.4 in Darve and Wooters. The algorithm is explained there.\\

Analysis: (*) please provide an example where we observe cubic convergence for bottom right side of matrices,
(*) return a metric such as ``number of QR iterations per singular value'', (*) time
your code in the square case for various $m=n$ and try to scale up to larger $m=n$
(say up to size $300$), 
compare with native julia functionality.\\

Rubrique: your code works for random matrices, is clearly written and the
analysis (stability, convergence and timing) is done, the choice made in the
design of the code are clearly explained in the comments of the code, you are able to clearly 
explain the code and answer questions 
during the oral presentation.

% The code needs to work for $m>n$

% How to get more points:

% (*) Write the code for complex matrix $A$ in input. The code needs to reduce
% to $A$ (which is complex ) to a real symmetric tridiagonal matrix.
% Then you can reuse your previous code for areal symmetric tridiagonal matrices.

% (*) Explore EAD: Early Aggressive Deflation

\end{document}
